% Answers to Chapter 7, translated from sv/back/facit/07.tex.

\facitavsnitt{c7}{Chapter 7}{Vectors}
\begin{facit}

\svar{1.1} \sv{a} $(7, 8)$ \sv{b} $(18, 22)$ \sv{c} $(-1, -4)$ \sv{d} $(-22, -28)$

\svar{1.2} \sv{a} $(0, 2)$ \sv{b} $(2, 6)$ \sv{c} $(-14, 16)$ \sv{d} $(18, 0)$

\svar{1.3} \sv{a} $\sqrt{5}$ \sv{b} $\sqrt{50} = 5\sqrt{2}$ \sv{c} $\sqrt{10}$
\sv{d} $\sqrt{125} = 5\sqrt{5}$

\svar{1.4} \sv{a} $x = 23$, $y = -5$ \sv{b} $x = -8$, $y = 2$

\svar{1.5} a) and d): $(6, -9) = 3\mathbf{u}$ and $(-12, 18) = -6\mathbf{u}$.

\svar{1.6} $a = 1.5$, $b = 3.5$

\svar{2.1} \sv{a} $5$ \sv{b} $5$ \sv{c} $5$ \sv{d} $8$ \sv{e} $4$

\svar{2.2} \sv{a} $\overrightarrow{AB} = (4, 3)$, $\abs{\overrightarrow{AB}} = 5$
\sv{b} $\overrightarrow{AB} = (3, -4)$, $\abs{\overrightarrow{AB}} = 5$
\sv{c} $\overrightarrow{BA} = (-4, -3)$: the same length, the opposite direction.
\sv{d} $\overrightarrow{BA} = -\overrightarrow{AB}$

\svar{2.3} \sv{a} $10$ \sv{b} $13$ \sv{c} $\sqrt{2} \approx 1.41$
\sv{d} $3\sqrt{2} \approx 4.24$ \sv{e} $7$

\svar{2.4} \sv{a} $45^{\circ}$, quadrant I, no correction
\sv{b} $135^{\circ}$, quadrant II, $+180^{\circ}$
\sv{c} $225^{\circ}$ (or $-135^{\circ}$), quadrant III, $+180^{\circ}$
\sv{d} $-45^{\circ}$ (or $315^{\circ}$), quadrant IV, no correction
\sv{e} $36.9^{\circ}$, quadrant I, no correction

\svar{2.5} \sv{a} $(8.66, 5.00)$ \sv{b} $(0, 5)$ \sv{c} $(-4, 0)$
\sv{d} $(-1.41, -1.41)$ \sv{e} $\sqrt{75.0 + 25} = \sqrt{100} = 10$

\svar{2.6} \sv{a} $11$ \sv{b} $-7$ \sv{c} $0$ \sv{d} $0$
\sv{e} The scalar product is zero, so the vectors are perpendicular.

\svar{2.7} \sv{a} Yes, the scalar product is $0$. \sv{b} No, the scalar product is $4$.
\sv{c} $k = 4$
\sv{d} For example $(2, 5)$, or a multiple of it.

\svar{2.8} \sv{a} $45^{\circ}$ \sv{b} $90^{\circ}$ \sv{c} About $81.9^{\circ}$
\sv{d} $180^{\circ}$ (opposite directions)

\svar{2.9} \sv{a} $(0.6, 0.8)$ \sv{b} $\sqrt{0.36 + 0.64} = 1$
\sv{c} $(-0.6, 0.8)$ \sv{d} $(9, 12)$ \sv{e} $(-9, -12)$

\svar{2.10} \sv{a} $(0, 0)$
\sv{b} The forces cancel each other out: the point is in equilibrium.
\sv{c} $(3, 7)$, with magnitude $\sqrt{58} \approx 7.62$
\sv{d} $(-3, -7)$, that is, exactly $\mathbf{F}_3$.

\svar{2.11} \sv{a} $(2, -3)$, $(-8, 12)$ and $(-2, 3)$
\sv{b} $(-8, 12)$ and $(-2, 3)$
\sv{c} $t = -6$
\sv{d} Yes: $\mathbf{u} = 1 \cdot \mathbf{u}$.

\svar{2.12} \sv{a} $(12, 5)\,\text{V}$ \sv{b} $13\,\text{V}$ \sv{c} About $22.6^{\circ}$
\sv{d} The voltages are perpendicular and add geometrically, like the two shorter sides of a
right-angled triangle: $\sqrt{12^2 + 5^2} = 13\,\text{V}$.

\svar{2.13} \sv{a} Both are $2$: the scalar product is commutative.
\sv{b} Both are $-2$: the scalar product is distributive over addition.
\sv{c} $\mathbf{u} \cdot \mathbf{u} = 5 = \abs{\mathbf{u}}^2$
\sv{d} $\mathbf{u} \cdot \mathbf{u} = \abs{\mathbf{u}}^2$

\svar{2.14} \sv{a} The length is given by Pythagoras' theorem, not by the sum:
$\abs{(3, 4)} = 5$.
\sv{b} The scalar multiplies both components: $(6, -2)$.
\sv{c} The scalar product is a number, not a vector: $8 + 15 = 23$.
\sv{d} The vector lies in quadrant II, so $180^{\circ}$ must be added: $126.9^{\circ}$.
\sv{e} A length is never negative: $\abs{-2\mathbf{u}} = 2\abs{\mathbf{u}}$.

\svar[Test yourself]{T} \sv{1} $\abs{\mathbf{u}} = 5$; the angle is about $-53.1^{\circ}$ (or
$306.9^{\circ}$), quadrant IV.
\sv{2} $\mathbf{w} = (5, -1)$

\end{facit}
